% ==============================================================================
% CAPITOLO 04 - MOTI RELATIVI E SISTEMI DI RIFERIMENTO NON INERZIALI
% ==============================================================================
\chapter{Moti Relativi e Sistemi di Riferimento Non Inerziali}
\label{chap:sistemi_non_inerziali}

\begin{tcolorbox}[enhanced, breakable, colback=navyblue!4!white, colframe=navyblue!80!black,
    coltitle=white, fonttitle=\bfseries\sffamily, title={Quadro Sinottico del Capitolo 04},
    sharp corners, rounded corners=south, boxrule=1pt]
\textbf{Collocazione Modulare:} Parte I -- Cinematica e Dinamica del Punto Materiale\\
\textbf{Manuale di Riferimento:} Focardi, Massa, Ugolini -- \emph{Fisica Generale: Meccanica e Termodinamica} (Capitolo 4)\\
\textbf{Obiettivi Didattici e Competenze:}\par
\begin{itemize}
    \item Formalizzare le trasformazioni cinematiche tra un sistema fisso (inerziale) $S$ e un sistema mobile $S'$ in moto rototraslatorio generico.
    \item Dimostrare rigorosamente il Teorema di Coriolis per la composizione delle velocità e delle accelerazioni ($\vec{a} = \vec{a}' + \vec{a}_{\text{trasc}} + \vec{a}_{\text{Cor}}$).
    \item Formulare il secondo principio nei sistemi non inerziali mediante l'introduzione delle forze apparenti (d'inerzia): forza di traslazione, forza centrifuga, forza di Eulero e forza di Coriolis.
    \item Analizzare il peso apparente in sistemi accelerati (ascensore e gravità artificiale).
    \item Studiare gli effetti della rotazione terrestre: variazione latitudinale di $\vec{g}_{\text{eff}}$, deviazione verso Est dei gravi in caduta libera e rotazione del piano del Pendolo di Foucault.
\end{itemize}
\end{tcolorbox}

\vspace{0.5em}

% ------------------------------------------------------------------------------
\section{Cinematica Relativa: Teorema di Composizione delle Accelerazioni}
\label{sec:cinematica_relativa}
% ------------------------------------------------------------------------------

Consideriamo due terne cartesiane ortonormali destre:
\begin{itemize}
    \item $S = (O, \ux, \uy, \uz)$: sistema di riferimento \textbf{fisso} (inerziale).
    \item $S' = (O', \ux', \uy', \uz')$: sistema di riferimento \textbf{mobile} (non inerziale), la cui origine $O'$ ha posizione $\vec{r}_{O'}(t)$ rispetto a $O$ e i cui assi ruotano con velocità angolare istantanea $\vec{\omega}(t)$.
\end{itemize}

\begin{center}
\begin{tikzpicture}[scale=1.2, >=Stealth]
    % Sistema Fisso S
    \coordinate (O) at (0,0);
    \draw[->, thick] (O) -- (2.5,0) node[below] {$x$};
    \draw[->, thick] (O) -- (0,2.5) node[left] {$y$};
    \filldraw[black] (O) circle (1.5pt) node[below left] {$O$ (Fisso $S$)};
    
    % Sistema Mobile S'
    \coordinate (Op) at (4,1.5);
    \begin{scope}[shift={(Op)}, rotate=25]
        \draw[->, thick, navyblue] (0,0) -- (2,0) node[below] {$x'$};
        \draw[->, thick, navyblue] (0,0) -- (0,2) node[left] {$y'$};
        \draw[->, ultra thick, purpleaccent] (0,0) -- (0, 1.2) node[above right] {$\vec{\omega}$};
    \end{scope}
    \filldraw[navyblue] (Op) circle (1.5pt) node[below] {$O'$ (Mobile $S'$)};
    
    % Vettore congiungente origini
    \draw[->, thick, forestgreen] (O) -- (Op) node[midway, below right] {$\vec{r}_{O'}$};
    
    % Punto P
    \coordinate (P) at (4.5, 4.2);
    \filldraw[crimsonred] (P) circle (2pt) node[above right] {$P$};
    
    % Vettori posizione
    \draw[->, very thick, crimsonred] (O) -- (P) node[midway, above left] {$\vec{r}$};
    \draw[->, very thick, navyblue] (Op) -- (P) node[midway, right] {$\vec{r}'$};
\end{tikzpicture}
\end{center}

\subsection{Posizione e Velocità Relativa}
Il vettore posizione assoluto $\vec{r}(t)$ del punto $P$ si esprime come somma vettoriale:
\begin{equation}
    \vec{r}(t) = \vec{r}_{O'}(t) + \vec{r}'(t) = \vec{r}_{O'}(t) + x'(t)\ux' + y'(t)\uy' + z'(t)\uz'
    \label{eq:pos_rel}
\end{equation}
Derivando rispetto al tempo nel sistema fisso $S$ e applicando le formule di Poisson per i versori mobili $\dv{\hat{\bm{u}}'_i}{t} = \vec{\omega} \times \hat{\bm{u}}'_i$:
\begin{align}
    \vec{v} &= \dv{\vec{r}}{t} = \dv{\vec{r}_{O'}}{t} + (\dot{x}'\ux' + \dot{y}'\uy' + \dot{z}'\uz') + \left(x'\dv{\ux'}{t} + y'\dv{\uy'}{t} + z'\dv{\uz'}{t}\right) \nonumber \\
    &= \vec{v}_{O'} + \vec{v}' + \vec{\omega}\times(x'\ux' + y'\uy' + z'\uz')
\end{align}

\begin{teorema}{Composizione delle Velocità (Galileo)}{teo:comp_velocita}
La velocità assoluta $\vec{v}$ è la somma della \textbf{velocità relativa} $\vec{v}'$ e della \textbf{velocità di trascinamento} $\vec{v}_{\text{trasc}}$:
\begin{equation}
    \vec{v} = \vec{v}' + \vec{v}_{\text{trasc}}
    \label{eq:comp_vel}
\end{equation}
dove la velocità di trascinamento rappresenta la velocità che avrebbe il punto $P$ se fosse solidale al sistema mobile $S'$:
\begin{equation}
    \vec{v}_{\text{trasc}} \coloneqq \vec{v}_{O'} + \vec{\omega}\times\vec{r}'
\end{equation}
\end{teorema}

\subsection{Teorema di Coriolis per le Accelerazioni}

Derivando nuovamente la velocità assoluta \eqref{eq:comp_vel} rispetto al tempo:
\begin{equation}
    \vec{a} = \dv{\vec{v}}{t} = \dv{\vec{v}_{O'}}{t} + \dv{\vec{v}'}{t} + \dv{\vec{\omega}}{t}\times\vec{r}' + \vec{\omega}\times\dv{\vec{r}'}{t}
\end{equation}
Calcolando la derivata temporale della velocità relativa $\vec{v}'$ vista da $S$:
\begin{equation}
    \dv{\vec{v}'}{t} = \vec{a}' + \vec{\omega}\times\vec{v}'
\end{equation}
e poiché $\dv{\vec{r}'}{t} = \vec{v}' + \vec{\omega}\times\vec{r}'$, sostituendo si ottiene:
\begin{equation}
    \vec{a} = \vec{a}' + \left[\vec{a}_{O'} + \dot{\vec{\omega}}\times\vec{r}' + \vec{\omega}\times(\vec{\omega}\times\vec{r}')\right] + 2\vec{\omega}\times\vec{v}'
\end{equation}

\begin{teorema}{Teorema di Coriolis delle Accelerazioni}{teo:coriolis_acc}
L'accelerazione assoluta $\vec{a}$ di un punto materiale si scompone in:
\begin{equation}
    \vec{a} = \vec{a}' + \vec{a}_{\text{trasc}} + \vec{a}_{\text{Cor}}
    \label{eq:teorema_coriolis}
\end{equation}
dove:
\begin{itemize}
    \item $\vec{a}' \coloneqq \ddot{x}'\ux' + \ddot{y}'\uy' + \ddot{z}'\uz'$ è l'\textbf{accelerazione relativa} (misurata dall'osservatore in $S'$).
    \item $\vec{a}_{\text{trasc}} \coloneqq \vec{a}_{O'} + \dot{\vec{\omega}}\times\vec{r}' + \vec{\omega}\times(\vec{\omega}\times\vec{r}')$ è l'\textbf{accelerazione di trascinamento}.
    \item $\vec{a}_{\text{Cor}} \coloneqq 2\vec{\omega}\times\vec{v}'$ è l'\textbf{accelerazione complementare (di Coriolis)}.
\end{itemize}
\end{teorema}

% ------------------------------------------------------------------------------
\section{Dinamica nei Sistemi Non Inerziali e Forze Apparenti}
\label{sec:forze_apparenti}
% ------------------------------------------------------------------------------

Nel sistema inerziale $S$, il secondo principio della dinamica stabilisce che $\vec{F}_{\text{reali}} = m\vec{a}$. Sostituendo la decomposizione di Coriolis \eqref{eq:teorema_coriolis}:
\begin{equation}
    \vec{F}_{\text{reali}} = m(\vec{a}' + \vec{a}_{\text{trasc}} + \vec{a}_{\text{Cor}}) \iff m\vec{a}' = \vec{F}_{\text{reali}} - m\vec{a}_{\text{trasc}} - m\vec{a}_{\text{Cor}}
\end{equation}

\begin{legge}{Equazione Fondamentale della Dinamica Relativa}{legge:dinamica_relativa}
Per un osservatore solidale a un riferimento non inerziale $S'$, il secondo principio assume formalmente la stessa struttura newtoniana:
\begin{equation}
    m\vec{a}' = \vec{F}_{\text{reali}} + \vec{F}_{\text{app}}
    \label{eq:dinamica_non_inerziale}
\end{equation}
a patto di aggiungere alla risultante delle interazioni fisiche reali $\vec{F}_{\text{reali}}$ le cosiddette \textbf{forze apparenti} (o fittizie/inerziali):
\begin{equation}
    \vec{F}_{\text{app}} \coloneqq \vec{F}_{\text{trasl}} + \vec{F}_{\text{cf}} + \vec{F}_{\text{Euler}} + \vec{F}_{\text{Cor}}
\end{equation}
\end{legge}

\begin{table}[htbp]
\centering
\caption{Quadro sinottico delle forze apparenti nei riferimenti non inerziali.}
\label{tab:forze_apparenti}
\begin{tabularx}{\textwidth}{p{3.2cm} p{3.8cm} p{3.2cm} X}
\toprule
\textbf{Forza Apparente} & \textbf{Espressione Analitica} & \textbf{Condizione} & \textbf{Effetto Fisico Principale} \\
\midrule
\textbf{Inerzia di Traslazione} & $\vec{F}_{\text{trasl}} = -m\vec{a}_{O'}$ & $\vec{a}_{O'} \neq \vec{0}$ & Spinta all'indietro in frenata o accelerazione rettilinea. \\
\textbf{Forza Centrifuga} & $\vec{F}_{\text{cf}} = -m\vec{\omega}\times(\vec{\omega}\times\vec{r}')$ & $\vec{\omega} \neq \vec{0}$ & Spinta radiale verso l'esterno dell'asse di rotazione. \\
\textbf{Forza di Eulero} & $\vec{F}_{\text{Euler}} = -m\dot{\vec{\omega}}\times\vec{r}'$ & $\dot{\vec{\omega}} \neq \vec{0}$ & Forza trasversa dovuta a variazioni di velocità angolare. \\
\textbf{Forza di Coriolis} & $\vec{F}_{\text{Cor}} = -2m\vec{\omega}\times\vec{v}'$ & $\vec{\omega} \neq \vec{0}, \, \vec{v}' \neq \vec{0}$ & Deviazione perpendicolare alla velocità relativa. \\
\bottomrule
\end{tabularx}
\end{table}

\begin{esame}[Natura delle Forze Apparenti]
Le forze apparenti non originano da interazioni fisiche con altri corpi (non hanno un corpo sorgente reale) e pertanto **non soddisfano il terzo principio di Newton** (non esiste alcuna forza di reazione corrispondente).
\end{esame}

% ------------------------------------------------------------------------------
\section{Applicazioni ed Effetti della Rotazione Terrestre}
\label{sec:applicazioni_terra}
% ------------------------------------------------------------------------------

La Terra ruota attorno al proprio asse polare con velocità angolare:
\begin{equation}
    \omega = \frac{2\pi}{T_{\text{siderale}}} = \frac{2\pi}{\SI{86164}{\second}} \approx \SI{7.292e-5}{\radian\per\second}
\end{equation}
Un laboratorio solidale alla superficie terrestre a latitudine $\lambda$ costituisce un sistema di riferimento non inerziale soggetto sia alla forza centrifuga che alla forza di Coriolis.

\subsection{Variazione Latitudinale della Gravità Effettiva}
Un corpo fermo sulla superficie ($\vec{v}' = \vec{0} \implies \vec{F}_{\text{Cor}} = \vec{0}$) è soggetto alla forza peso gravitazionale pura $\vec{P}_0 = m\vec{g}_0$ e alla forza centrifuga $\vec{F}_{\text{cf}} = m\omega^2 R_T \cos\lambda\,\hat{\bm{u}}_\perp$:
\begin{equation}
    \vec{g}_{\text{eff}} = \vec{g}_0 + \omega^2 R_T\cos\lambda\,\hat{\bm{u}}_\perp
\end{equation}
\begin{itemize}
    \item \textbf{Ai Poli ($\lambda = \pm 90^\circ$)}: $\cos\lambda = 0 \implies g_{\text{eff}} = g_0 \approx \SI{9.832}{\metre\per\second\squared}$.
    \item \textbf{All'Equatore ($\lambda = 0^\circ$)}: la forza centrifuga è massima e opposta a $\vec{g}_0$:
    \begin{equation}
        g_{\text{eff}} = g_0 - \omega^2 R_T \approx 9{,}832 - (7{,}292\times 10^{-5})^2 \times (6{,}371\times 10^6) \approx \SI{9.780}{\metre\per\second\squared}
    \end{equation}
    \item \textbf{Deviazione del Filo a Piombo}: a latitudini intermedie, la direzione della verticale apparente (perpendicolare alla superficie degli oceani) devia verso l'equatore rispetto alla direzione radiale verso il centro della Terra di un angolo massimo $\delta \approx 0{,}1^\circ$ a $\lambda = 45^\circ$.
\end{itemize}

\subsection{Deviazione verso Est dei Gravi in Caduta Libera}
Un grave lasciato cadere da fermo da un'altezza $h$ cade verso il basso con velocità relativa diretta lungo la verticale $\vec{v}' \approx -g t\uz'$. La forza di Coriolis agente su di esso vale:
\begin{equation}
    \vec{F}_{\text{Cor}} = -2m\vec{\omega}\times(-gt\uz') = 2mg t (\vec{\omega}\times\uz')
\end{equation}
A latitudine $\lambda$, il prodotto vettoriale genera una spinta diretta verso **Est** (verso entrante):
\begin{equation}
    a_x' = 2\omega g t \cos\lambda
\end{equation}
Integrando due volte nel tempo da $t=0$ al tempo di caduta $t_{\text{cad}} = \sqrt{\frac{2h}{g}}$:
\begin{equation}
    x_{\text{Est}} = \int_0^{t_{\text{cad}}} \left(\int_0^t 2\omega g \tau \cos\lambda\,\dd\tau\right)\dd t = \frac{1}{3}\omega g \cos\lambda\,t_{\text{cad}}^3 = \mathbf{\frac{2\sqrt{2}}{3}\omega\cos\lambda \sqrt{\frac{h^3}{g}}}
    \label{eq:deviazione_est}
\end{equation}

\subsection{Pendolo di Foucault}
Un pendolo sferico di grande lunghezza oscilla liberamente su un piano. Per effetto della componente verticale della velocità angolare terrestre $\omega_z = \omega\sin\lambda$, la forza di Coriolis deflette la massa sempre verso destra (nell'emisfero boreale), facendo ruotare il piano di oscillazione in senso orario con periodo:
\begin{equation}
    T_F = \frac{24\text{ ore}}{\sin\lambda}
\end{equation}
Ai poli ($\lambda = 90^\circ$), $T_F = 24\text{ h}$ (la Terra ruota sotto il pendolo), mentre all'equatore ($\lambda = 0$) il piano di oscillazione non ruota ($T_F \to \infty$).

% ------------------------------------------------------------------------------
\section{Esercizi Risolti ed Applicativi con Analisi delle Forze Apparenti}
\label{sec:esercizi_non_inerziali}
% ------------------------------------------------------------------------------

\begin{esercizio}[Pallina su Asta Rotante e Diagramma nel Riferimento Non Inerziale]
\textbf{Testo:}
Un anellino di massa $m$ può scorrere senza attrito lungo un'asta orizzontale rettilinea posta in rotazione attorno a un asse verticale passante per $O$ con velocità angolare costante $\omega$. All'istante $t=0$, l'anellino si trova a distanza $r_0 > 0$ dal centro con velocità radiale nulla $\dot{r}(0) = 0$. Determinare la legge oraria $r(t)$ e analizzare il diagramma delle forze nel riferimento rotante $S'$.
\end{esercizio}

\begin{center}
\begin{tikzpicture}[scale=1.2, >=Stealth]
    % Asse di rotazione e asta
    \draw[thick, dashed] (0,-1) -- (0,2) node[above] {Asse di rotazione $\vec{\omega} = \omega\uz$};
    \draw[very thick, brown!80!black] (0,0) -- (5, 0) node[right] {Asta guida};
    \filldraw[black] (0,0) circle (2pt) node[below left] {$O$};
    
    % Anello di massa m
    \coordinate (M) at (3,0);
    \draw[ultra thick, fill=blue!20, draw=navyblue] (2.7, -0.3) rectangle (3.3, 0.3);
    \filldraw[black] (M) circle (1.5pt) node[below=12pt] {$m$};
    
    % Vettori forza nel piano orizzontale (riferimento S')
    \draw[->, ultra thick, crimsonred] (M) -- (4.8, 0) node[above] {$\vec{F}_{\text{cf}} = m\omega^2 r\,\ur'$};
    \draw[->, ultra thick, purpleaccent] (M) -- (3, -1.5) node[right] {$\vec{F}_{\text{Cor}} = -2m\omega\dot{r}\,\utheta'$};
    \draw[->, ultra thick, forestgreen] (M) -- (3, 1.5) node[right] {$\vec{N}_\theta$ (Reazione guida)};
    
    % Velocità radiale
    \draw[->, very thick, navyblue] (M) -- (4.0, 0) node[below] {$\dot{r}\,\ur'$};
\end{tikzpicture}
\end{center}

\begin{tcolorbox}[enhanced, breakable, colback=forestgreen!5!white, colframe=forestgreen!80!black,
    title={\textbf{Analisi delle Forze nel Riferimento Rotante $S'$}}, fonttitle=\bfseries\sffamily]
\begin{itemize}
    \item \textbf{Lungo la direzione radiale $\ur'$:}
    \begin{itemize}
        \item \textbf{Forza Centrifuga ($\vec{F}_{\text{cf}} = m\omega^2 r\,\ur'$):} è una forza apparente diretta radialmente verso l'esterno; è l'unica forza agente lungo l'asta ed è responsabile dell'accelerazione radiale $\ddot{r}$.
        \item \textbf{Nessuna forza reale radiale ($F_{r,\text{reale}} = 0$):} l'asta è perfettamente liscia, quindi non può esercitare forze di attrito tangenziali al moto radiale.
    \end{itemize}
    \item \textbf{Lungo la direzione trasversa $\utheta'$:}
    \begin{itemize}
        \item \textbf{Forza di Coriolis ($\vec{F}_{\text{Cor}} = -2m\vec{\omega}\times\vec{v}' = -2m\omega\dot{r}\utheta'$):} spinge l'anellino contro la parete laterale dell'asta (nel verso opposto alla rotazione).
        \item \textbf{Reazione Normale Trasversa ($\vec{N}_\theta$):} la parete dell'asta bilancia esattamente la forza di Coriolis: $\vec{N}_\theta = -\vec{F}_{\text{Cor}} = +2m\omega\dot{r}\utheta'$.
    \end{itemize}
    \item \textbf{Lungo la verticale $\uz$:} la reazione normale verticale bilancia la forza peso ($\vec{N}_z + \vec{P} = \vec{0}$).
\end{itemize}
\end{tcolorbox}

\textbf{Risoluzione Analitica:}
L'equazione del moto radiale è:
\begin{equation}
    m\ddot{r} = m\omega^2 r \implies \ddot{r}(t) - \omega^2 r(t) = 0
\end{equation}
Con le condizioni iniziali $r(0) = r_0$ e $\dot{r}(0) = 0$, la soluzione è:
\begin{equation}
    \mathbf{r(t) = r_0 \cosh(\omega t)}
\end{equation}

\vspace{0.8em}

\begin{esercizio}[Ascensore Accelerato e Peso Apparente]
\textbf{Testo:}
Una persona di massa $m = \SI{70}{\kilogram}$ si trova su una bilancia all'interno di un ascensore che accelera verso l'alto con modulo $a_0 = \SI{2.0}{\metre\per\second\squared}$. Determinare la lettura della bilancia (peso apparente) e confrontare l'analisi nei sistemi inerziale e non inerziale.
\end{esercizio}

\begin{center}
\begin{tikzpicture}[scale=1.1, >=Stealth]
    % Cabina ascensore
    \draw[very thick] (-2,-1) rectangle (2, 2.5);
    \node[above] at (0, 2.5) {Cabina $S'$};
    \draw[->, ultra thick, purpleaccent] (2.5, 0.5) -- (2.5, 1.8) node[right] {$\vec{a}_0 = a_0\uy$};
    
    % Bilancia e persona
    \filldraw[fill=gray!40] (-1.2, -1) rectangle (1.2, -0.7);
    \node at (0, -0.85) {\small Bilancia};
    \filldraw[fill=blue!10, draw=navyblue, thick] (-0.6, -0.7) rectangle (0.6, 0.8);
    \node at (0, 0.05) {$m$};
    
    % Forze
    \draw[->, ultra thick, forestgreen] (0, -0.7) -- (0, 1.3) node[right] {$\vec{N}$ (Lettura bilancia)};
    \draw[->, ultra thick, navyblue] (0, 0.05) -- (0, -1.5) node[right] {$\vec{P} = m\vec{g}$};
    \draw[->, ultra thick, crimsonred] (-0.3, 0.05) -- (-0.3, -1.0) node[left] {$\vec{F}_{\text{in}} = -m\vec{a}_0$};
\end{tikzpicture}
\end{center}

\begin{tcolorbox}[enhanced, breakable, colback=forestgreen!5!white, colframe=forestgreen!80!black,
    title={\textbf{Confronto tra Sistema Inerziale e Riferimento Non Inerziale}}, fonttitle=\bfseries\sffamily]
\begin{itemize}
    \item \textbf{Osservatore Inerziale a Terra ($S$):}
    Vede la persona accelerare verso l'alto con $\vec{a} = \vec{a}_0$. La seconda legge di Newton stabilisce che la risultante delle forze reali fornisce l'accelerazione:
    \begin{equation}
        N - mg = m a_0 \implies N = m(g + a_0) = 70 \times (9{,}81 + 2{,}0) = \mathbf{\SI{826.7}{\newton}}
    \end{equation}
    \item \textbf{Osservatore Non Inerziale nell'Ascensore ($S'$):}
    Vede la persona in quiete ($\vec{a}' = \vec{0}$). All'equilibrio statico agiscono il peso reale $\vec{P} = -mg\uy$, la forza apparente d'inerzia $\vec{F}_{\text{in}} = -m\vec{a}_0 = -m a_0\uy$ e la reazione normale $\vec{N}$:
    \begin{equation}
        \sum \vec{F}' = \vec{0} \implies N - mg - m a_0 = 0 \implies \mathbf{N = m(g + a_0)}
    \end{equation}
    La bilancia misura la forza normale $N$, pari al \textbf{peso apparente} $P_{\text{app}} = m g_{\text{eff}} > mg$.
\end{itemize}
\end{tcolorbox}
